\abstract
Dependency-Injection is a technique used to dynamically add and load dependencies within an application.
This can be wanted but also unwanted.

This research report will evaluate rather, enabling Dependency-Injection in Java comes with high risk and poses a high threat to the application, or, if implementing it properly is actually worth the risk, comes with benefits and if it is actually possible to prevent unwanted injections and scenarios to happen.
For evaluation purpose we will conduct an experiment, showcasing the worst-case-scenario, and evaluate based on that what can be done prevent this.
Additionally, different types of injecting ("attacking") will be discussed.

This report is written by \textit{Lukas Weber} and \textit{Gianluca Biancolin} as part of the \textbf{SEAR}\footnote{Security and Research} module from the FHTenL\footnote{Fontys University of Applied Science / Technic and Logistic, Venlo, the Netherlands}.
The report will be graded.

\clearpage